\documentclass[11pt,a4paper]{article}
\usepackage[utf8]{inputenc}
\usepackage[T1]{fontenc}
\usepackage[french]{babel}
\usepackage[margin=2.2cm]{geometry}
\usepackage{booktabs}
\usepackage{array}
\usepackage{xcolor}
\usepackage{amssymb}
\usepackage{hyperref}
\usepackage{fancyhdr}
\usepackage{verbatim}

\definecolor{accent}{HTML}{2D6BE8}
\definecolor{warn}{HTML}{C45C26}

\hypersetup{colorlinks=true, linkcolor=accent, urlcolor=accent}

\setlength{\headheight}{14pt}
\pagestyle{fancy}
\fancyhf{}
\lhead{Corrigé Vision Desk --- Computer Vision Finance M2}
\cfoot{\thepage}

\title{%
  \textbf{Corrigé : Vision Desk}\\[0.35em]
  \large ViT, JSON et MCP grounded\\[0.5em]
  \normalsize Computer Vision in Finance --- Master 2 Finance, Data et IA
}
\author{Dhiya Lecheheb}
\date{2026}

\begin{document}
\maketitle
\thispagestyle{empty}

\section{Lecture du corrigé}

Ce document répond à la \textbf{tâche fermée} du cas avec un protocole de
référence et des traces types. La section~\ref{sec:lab-run} consigne les métriques
d'une exécution du notebook (\texttt{SEED=42}). Les autres F1 restent pédagogiques.
Non négociable : ViT (pas LSTM), schéma JSON avant le test, quatre tools MCP,
aucune affirmation inventée.

\section{Labo de référence (Colab)}

Notebook exécutable, aligné sur ce corrigé :

\begin{center}
\renewcommand{\arraystretch}{1.15}
\begin{tabular}{@{}p{3.2cm}p{11.4cm}@{}}
\toprule
\textbf{Ressource} & \textbf{Lien} \\
\midrule
Dépôt & \href{https://github.com/Amzil-AI/Vision-desk}{github.com/Amzil-AI/Vision-desk} \\
Notebook & \texttt{vision\_desk\_colab.ipynb} \\
Ouvrir dans Colab & \href{https://colab.research.google.com/github/Amzil-AI/Vision-desk/blob/main/vision_desk_colab.ipynb}{colab.research.google.com/.../vision\_desk\_colab.ipynb} \\
\bottomrule
\end{tabular}
\end{center}

Le notebook couvre : images finance synthétiques, split par document, fine-tune
de tête ViT vs baseline majorité, contrat JSON (\texttt{vision.extract}), quatre
stubs MCP, plus traces grounded et refus. Les sorties de référence sont dans
\texttt{solution/lab-reference/} (ci-dessous). Les groupes peuvent remplacer les images
synthétiques par leurs captures si le protocole reste identique.

\subsection{Exécution enregistrée (notebook, \texttt{SEED=42})}
\label{sec:lab-run}

Le code du notebook Colab a été exécuté de bout en bout (CPU local, octobre 2026 ;
même logique et seed que sur
\href{https://colab.research.google.com/github/Amzil-AI/Vision-desk/blob/main/vision_desk_colab.ipynb}{Colab}).
Les artefacts sont copiés dans le dépôt sous \texttt{solution/lab-reference/}
(\texttt{metrics.json}, \texttt{grounded\_trace.json}, \texttt{refusal\_trace.json}).

\begin{center}
\renewcommand{\arraystretch}{1.15}
\begin{tabular}{@{}p{4.2cm}p{10.4cm}@{}}
\toprule
\textbf{Métrique} & \textbf{Valeur (documents held-out)} \\
\midrule
Images train / val & 96 / 32 (8 documents val) \\
F1 macro baseline majorité & 0{,}056 \\
F1 macro ViT-Tiny (tête, 8 epochs) & 1{,}000 \\
Bat la baseline ? & oui \\
Fichiers & \texttt{metrics.json}, \texttt{grounded\_trace.json}, \texttt{refusal\_trace.json} \\
\bottomrule
\end{tabular}
\end{center}

{\color{warn}\textbf{Lecture.}
Un F1 macro $=1{,}0$ sur les templates \emph{synthétiques} est attendu : layouts
propres et classes visuellement distinctes. Des captures Bloomberg / PDF réelles
donneront un F1 plus bas et une matrice de confusion plus riche --- un résultat
négatif honnête est acceptable si le split par document et les traces MCP sont
corrects.}

\paragraph{Exemple \texttt{vision.extract} (val, avec ticker).}
\begin{verbatim}
{
  "type": "candlestick_chart",
  "confidence": 0.989,
  "ticker": "AAPL",
  "fields": {"title": "candlestick_chart__doc05_p0", "period": "2024"},
  "source_tool": "vision.extract"
}
\end{verbatim}

\paragraph{Réponse grounded (issue de \texttt{lab-reference/grounded\_trace.json}).}
« \texttt{vision.classify} classe en \texttt{candlestick\_chart} (score 0{,}989).
\texttt{vision.extract} renvoie \texttt{ticker=AAPL}. \texttt{market.lookup} :
Apple Inc.\ (USD). Je n'affirme aucun prix de clôture : aucun champ prix dans le JSON. »

\paragraph{Refus (issue de \texttt{lab-reference/refusal\_trace.json}).}
« Je ne peux pas confirmer de ticker : extraction null ou \texttt{market.lookup}
en échec. Je refuse d'inventer un symbole. Relancer avec un crop plus large / autre DPI. »

\section{Réponse à la tâche fermée (structure)}

\begin{center}
\renewcommand{\arraystretch}{1.2}
\begin{tabular}{@{}p{4.0cm}p{10.6cm}@{}}
\toprule
\textbf{Exigence} & \textbf{Ce qu'un bon groupe montre} \\
\midrule
1. Assertion vision & Classe + score via \texttt{vision.classify} ; champs via \texttt{vision.extract} \\
2. Justification outil & Chaque phrase cite un tool call et un champ JSON \\
3. Refus & Si ticker null / tool KO $\rightarrow$ refus ou relance, jamais inventer \\
\bottomrule
\end{tabular}
\end{center}

\section{Protocole retenu}

\begin{center}
\renewcommand{\arraystretch}{1.15}
\begin{tabular}{@{}p{3.8cm}p{10.8cm}@{}}
\toprule
\textbf{Choix} & \textbf{Valeur} \\
\midrule
Entrée & Captures de pages / charts, 2--4 classes \\
Split & Par document (labo : 75\,\% / 25\,\% docs held-out ; 96 / 32 images) \\
Modèle & ViT-Tiny (ou Small) préentraîné + tête fine-tunée \\
Baseline & Classe majoritaire, ou logits sur pixels réduits (séance~1) \\
JSON & \texttt{type}, \texttt{confidence}, \texttt{ticker?} (nullable), \texttt{fields} \\
MCP & \texttt{vision.classify}, \texttt{vision.extract}, \texttt{market.lookup}, \texttt{knowledge.search} \\
LLM & N'affirme un fait que s'il apparaît dans un JSON d'outil \\
\bottomrule
\end{tabular}
\end{center}

\section{Schéma JSON (contrat)}

\begin{verbatim}
{
  "type": "candlestick_chart",
  "confidence": 0.989,
  "ticker": "AAPL",
  "fields": {
    "title": "candlestick_chart__doc05_p0",
    "period": "2024"
  },
  "source_tool": "vision.extract"
}
\end{verbatim}

Même payload que l'exemple labo ci-dessus (title = stem de l'image).
\texttt{ticker} peut être \texttt{null}. Un JSON invalide (champ requis manquant)
est un échec pipeline même si la classe est correcte.

\section{Résultats vision (ordre de grandeur attendu)}

\begin{center}
\renewcommand{\arraystretch}{1.2}
\begin{tabular}{@{}lcc@{}}
\toprule
\textbf{Modèle} & \textbf{F1 macro (held-out)} & \textbf{Lecture} \\
\midrule
Labo de réf. (synth.) & maj.\ 0{,}06 / ViT 1{,}00 & Voir \S\ref{sec:lab-run} ; pas une promesse prod. \\
Baseline majorité (données réelles) & $\approx$ 0{,}25--0{,}40 & Selon déséquilibre \\
ViT fine-tune & nettement plus haut & Si layout stable \\
ViT sur nouveau template & souvent en baisse & Domain shift attendu \\
\bottomrule
\end{tabular}
\end{center}

{\color{warn}\textbf{Lecture.}
Si le held-out mélange des pages du même PDF que le train, le F1 est gonflé
(fuite documentaire). Analogue de \texttt{train\_test\_split} aléatoire dans Signal Desk.}

\section{Trace grounded (exigences 1 + 2)}

\paragraph{Question utilisateur.} « Quel type d'image est-ce, et quel ticker apparaît\,? »

\paragraph{Tools.}
\begin{enumerate}
  \item \texttt{vision.classify} $\rightarrow$ \texttt{\{class: candlestick\_chart, score: 0.989\}}
  \item \texttt{vision.extract} $\rightarrow$ JSON ci-dessus (\texttt{ticker: "AAPL"})
  \item \texttt{market.lookup("AAPL")} $\rightarrow$ métadonnées (nom, devise) --- stub OK
  \item \texttt{knowledge.search("candlestick")} $\rightarrow$ glossaire court --- stub OK
\end{enumerate}

\paragraph{Réponse finale acceptable.}
« \texttt{vision.classify} classe l'image comme graphique en chandeliers (0{,}989).
\texttt{vision.extract} renvoie \texttt{ticker=AAPL}.
\texttt{market.lookup} confirme les métadonnées Apple Inc.
Je n'affirme aucun prix de clôture : aucun champ prix n'est présent dans le JSON. »

\section{Refus après échec d'outil (exigence 3)}

\paragraph{Situation.} \texttt{vision.extract} renvoie \texttt{ticker: null}
(ou \texttt{market.lookup} renvoie une erreur).

\paragraph{Acceptable.}
« Je ne peux pas confirmer de ticker : l'extraction renvoie null.
Je refuse d'inventer un symbole. Relancer avec un crop plus large / autre DPI. »

\paragraph{Inacceptable.}
« C'est clairement Tesla. » --- aucune source tool.

\section{Anti-patterns sanctionnés}

\begin{itemize}
  \item LSTM ou Conv1D sur l'image « parce que deep learning »\,;
  \item affirmation LLM sans tool call\,;
  \item appeler \texttt{vision.*} puis ignorer le JSON\,;
  \item présenter un glossaire RAG comme une extraction OCR\,;
  \item tools d'écriture / mutation dans le labo (MCP en lecture seule uniquement).
\end{itemize}

\section{Limites types (à recopier dans la note 4 pages)}

\begin{enumerate}
  \item \textbf{Layout shift} --- nouveau template $\Rightarrow$ chute de F1.
  \item \textbf{Capture} --- DPI, barre d'UI, compression JPEG.
  \item \textbf{Hallucination} --- le LLM complète ce que le JSON n'a pas.
\end{enumerate}

\section{Lien avec Signal Desk}

\begin{center}
\renewcommand{\arraystretch}{1.15}
\begin{tabular}{@{}p{5.5cm}p{8.5cm}@{}}
\toprule
\textbf{Signal Desk} & \textbf{Vision Desk} \\
\midrule
Série OHLCV + LSTM & Image + ViT \\
Baseline naïve / coûts & Baseline classe + grounding MCP \\
Walk-forward & Split par document + held-out templates \\
Signal mort après coûts = OK & Refus après tool KO = OK \\
\bottomrule
\end{tabular}
\end{center}

\section{Ressources vidéo}

\begin{itemize}
  \item \href{https://www.youtube.com/watch?v=eMlx5fFNoYc}{3Blue1Brown --- Attention}\,;
  \item \href{https://www.youtube.com/watch?v=TrdevFK_am4}{Kilcher --- ViT paper}\,;
  \item \href{https://www.youtube.com/watch?v=ovB0ddFtzzA}{Vision Transformer in PyTorch}\,;
  \item \href{https://www.youtube.com/watch?v=CQywdSdi5iA}{MCP --- Anthropic}\,;
  \item \href{https://www.youtube.com/watch?v=LYfr7qusVSs}{MCP beginners Python}.
\end{itemize}

\end{document}
